\SourceLeaf{45}
\[ \cos a \cos b = \tfrac12\cos(a-b)+\tfrac12\cos(a+b). \]

\par\noindent\hspace*{\parindent}Now take $\cos\varphi$, which, multiplied by $\cos\varphi$, gives $\cos\varphi . \cos\varphi$. This product is represented by $\cos\varphi^2$ if we actually multiply, or by $\tfrac12+\tfrac12\cos 2\varphi$ if we resolve it into a sum by the formula just stated. Therefore $\cos\varphi^2 = \tfrac12+\tfrac12\cos 2\varphi$. Proceeding in the same manner to higher powers, we gradually obtain the following formulas

\[ \cos\varphi^3 = \tfrac34\cos\varphi + \tfrac14\cos 3\varphi ; \]
\[ \cos\varphi^4 = \tfrac38 + \tfrac48\cos 2\varphi + \tfrac18\cos 4\varphi ; \]
\[ \cos\varphi^5 = \tfrac{10}{16}\cos\varphi + \tfrac{5}{16}\cos 3\varphi + \tfrac{1}{16}\cos 5\varphi ; \]
\[ \cos\varphi^6 = \tfrac{10}{32} + \tfrac{15}{32}\cos 2\varphi + \tfrac{6}{32}\cos 4\varphi + \tfrac{1}{32}\cos 6\varphi ; \]
\[ \cos\varphi^7 = \tfrac{35}{64}\cos\varphi + \tfrac{21}{64}\cos 3\varphi + \tfrac{7}{64}\cos 5\varphi + \tfrac{1}{64}\cos 7\varphi ; \]
\[ \cos\varphi^8 = \tfrac{35}{128} + \tfrac{56}{128}\cos 2\varphi + \tfrac{28}{128}\cos 4\varphi + \tfrac{8}{128}\cos 6\varphi + \tfrac{1}{128}\cos 8\varphi . \]

\par\noindent The general law is readily derived from these calculations. In even powers of $\cos\varphi$ only cosines of even multiples of $\varphi$ appear, the greatest multiple being determined by the given exponent. In odd powers, by the same rule, cosines of odd multiples appear. The even powers, however, do not follow exactly the same law as the odd powers.

$\alpha$. The initial term in the even powers, which is multiplied by $\cos 0\varphi = 1$, has a coefficient whose denominator, like that of all the others, is a power of 2 with exponent one less than the exponent of the proposed power. The numerator is in substance a binomial coefficient belonging to the proposed power, specifically half of the middle one when the coefficients are set out in order. The other terms have a multiple of $\varphi$ equal to twice the index of the term. In their coefficients they have the same denominator as the initial term, and for numerator a binomial coefficient of the same power, as many places after the one in the initial term as the index of the term indicates. This law can be expressed generally as follows. The series for $\cos\varphi^{2m}$ has the initial term\HistoricalNote{*)}{*)}{\textsc{Euler}'s expression of this law, which is not reduced to binomial coefficients, agrees in substance with the proposed expression but seems less concise.

I hardly need to mention that I always denote binomial coefficients in the manner of \textsc{Hindenburg}, and that, following the distinguished\ \textsc{Kaestner}, I always distinguish the initial term of a series from the subsequent terms.}

\[ \frac{1}{2^{2m-1}} . \frac{{}^{2m}\mathfrak{B}^{(m)}}{2}, \]
